%%%%%%%%%%%%%%%%%%%%%%%%%%%%%%%%%%%%%%%%%%%%%%%%%%%%%%%%%%%%%%%%%%%%%%%%%%%%%%%%
%%%%%%%%%%%%%%%%%%%%%%%%%%%%%%%%%%%%%%%%%%%%%%%%%%%%%%%%%%%%%%%%%%%%%%%%%%%%%%%%
% INTRODUCCIÓN %
% Explicar el contexto del proyecto: pq lo has hecho, cual es la motivacion, q esperas conseguir % 
%%%%%%%%%%%%%%%%%%%%%%%%%%%%%%%%%%%%%%%%%%%%%%%%%%%%%%%%%%%%%%%%%%%%%%%%%%%%%%%%


\cleardoublepage
\chapter{Introducción}
\label{chap:introduccion}
\pagenumbering{arabic}

En estos útlimos años, los modelos de lenguaje (LLMs) están cambian nuestra relación con la tecnología. Herramientas como ChatGPT o Gemini que aparecieron hace un tiempo en la industria, o herramientas mas actuales como Claude con gran potencial y variedad de uso han demostrado que en la mayoría de casos, hacer un buen uso de ellas puede ayudarte a hacer tareas que tradicionalmente requerían conocimientos técnicos avanzados.

Este cambio ha creado la necesidad de actualizarse y aprender el manejo adecuado de estas herramientas, porque con un buen uso y manejo de ellas se abren nuevas posibilidades en diferentes campos, como la redaccion de textos, la ayuda en generacion de codigo, o como en este caso la generacion automatica de codigo para proyectos.


\section{Contexto}
\label{sec:contexto}

En la actualidad, el desarrollo de webs es una tarea altamente demandada por todo tipo de sectores.
Sin embargo, construir una aplicación desde cero requiere conocimientos técnicos sobre varios ambitos: lenguajes de programación, bases de datos, frameworks, protocolos de comunicación, despliegue en producción, seguridad y diseño visual, entre otras. Esto implica que exiten un gran número de usuarios que tienen ideas y datos interesantes pero no tienen los conocimientos necesarios para llevarlas a una aplicación funcional.

Además, hay gran cantidad de datos accesibles públicamente a través de APIs, ofreciendo datos sobre gran variedad de temas: catálogos de productos, datos meteorológicos, criptomonedas, información geográfica, deportiva, etc. La mayoría de estos datos están perfectamente estructurados y disponibles de forma gratuita, pero requieren conocimientos técnicos para ser parseados, entendidos y presentados de forma útil al usuario final.

WebBuilder ofrece una vía accesible para que cualquier usuario, independientemente de su nivel técnico, pueda transformar una fuente de datos pública en un sitio web funcional y desplegable. 


\section{Motivación}
\label{sec:motivacion}

La motivacion para realizar este proyecto surge del interés en tres áreas diferentes, pero capaces de juntarse en un único proyecto.

En primer lugar, la experiencia previa con Django, framework utilizado en la asignatura de Servicios Telemáticos y en el cual me sentí muy cómodo. La facilidad de comienzo de una nueva aplicacion web, mas la cantidad de funcionalidades y personalizacion que ofrece, junto con la facilidad de organización de un proyecto grande hizo que fuese uno de los requisitos para empezar con esta base.

En segundo lugar, la inteligencia artificial. En este campo no me refiero a utilizarla como usuario final, sino a entender cómo funciona por dentro y aprender a controlar estas herramientas de forma adecuada. Un uso superficial de los modelos de lenguaje, sin comprender qué está ocurriendo realmente, genera una dependencia que puede volverse en contra del desarrollador. Aprender prompt engineering correctamente, saber interpretar las respuestas del modelo y conocer sus limitaciones reales permite avanzar mucho más rápido y con criterio, siempre teniendo conocimientos sólidos sobre el problema que se quiere resolver.

En tercer lugar, otro área en la que estaba interesado es la automatizacion de procesos, queriendo aprender los conocimientos necesarios para poder crear flujos de trabajo automaticos. En este caso, crear procesos automaticos hace que todo sea mas sencillo y rapido, pudiendo no solo lanzar proyectos con un simple click, sino tambien poder entender de forma actualizada como esta funcionando tu proyecto por dentro con datos y métricas reales, teniendo de esta forma control real sobre que sucede.


\section{Estructura de la memoria}
\label{sec:estructura}

Esta memoria se organiza en siete capítulos, cada uno enfocado en un aspecto diferente del proyecto:

\begin{itemize}
    \item \textbf{Capítulo 2 — Objetivos:} recoge los objetivos generales y específicos del proyecto, junto con la planificación temporal seguida durante el desarrollo y el resumen de las reuniones mantenidas con el tutor.

    \item \textbf{Capítulo 3 — Estado del arte:} describe las tecnologías, herramientas y conceptos que han hecho posible la construcción de WebBuilder: el framework web, los modelos de lenguaje, las herramientas de automatización e infraestructura y las tecnologías de frontend.

    \item \textbf{Capítulo 4 — Diseño e implementación:} aborda el desarrollo del sistema por sprints desde la definición del proyecto hasta el cierre, con especial atención a las decisiones técnicas tomadas en cada sprint, los problemas encontrados y la evolución visual de la plataforma.

    \item \textbf{Capítulo 5 — Experimentos y validación:} presenta los casos de uso reales generados con la plataforma, el análisis de los errores detectados durante el desarrollo y las métricas del sistema obtenidas tras varias semanas de uso.

    \item \textbf{Capítulo 6 — Resultados:} recoge los resultados obtenidos, evaluando en qué medida el sistema cumple con los requisitos planteados y mostrando los sitios generados como evidencia del funcionamiento real de la plataforma.

    \item \textbf{Capítulo 7 — Conclusiones:} cierra la memoria con la consecución de objetivos, la aplicación de los conocimientos del grado, las lecciones aprendidas durante el desarrollo y las posibles líneas de trabajo futuro.
\end{itemize}